\clearpage
\thispagestyle{empty}
\phantomsection
\addcontentsline{toc}{chapter}{Aynı Yazardan: Bilge ve Yonga}
{\Large\bfseries\color{bolumrenk} Aynı Yazardan: Bilge ve Yonga}
\vspace{0.5cm}
\noindent\rule{\textwidth}{2pt}
\vspace{0.8cm}

\noindent
Bilgisayar neden ısınır, işlemci bir buyruğu nasıl anlar, sayılar bir yongada nasıl saklanır? Bu kitabın anlattığı kavramların çoğu aslında çocukluk merakının doğal uzantısıdır. \textbf{Bilge ve Yonga}, bu soruların yanıtlarını okuma bilen çocuklara hikayelerle anlatan, ücretsiz ve açık erişimli bir resimli kitap serisidir. Bilge sekiz yaşında meraklı bir kız çocuğu, Yonga ise bir yongadan doğmuş, havada süzülen sevimli bir robottur.

\medskip
\noindent
Seri dört ciltten oluşur ve her cilt kendi içinde bir konuyu baştan sona işler:

\medskip
\begin{itemize}\setlength{\itemsep}{2pt}
    \item \textbf{Kumdan Bilgisayara} (7--10 yaş): bir avuç kumdan yongaya, yongadan çalışan bir bilgisayara uzanan yolculuk.
    \item \textbf{Hız ve Güç} (7--10 yaş): bilgisayarı hızlı ve güvenilir kılan nedir, başarımın sırları nelerdir.
    \item \textbf{Buyrukların Dünyası} (9--12 yaş): bilgisayara ne yapacağını nasıl söyleriz, buyrukların dili nasıl kurulur.
    \item \textbf{İşlemcinin İçi} (9--12 yaş): işlemcinin kapağını açmak; toplayıcılar, eldeler ve işaret taşıyan parçalar.
\end{itemize}

\medskip
\noindent
Elliyi aşkın kitabın tamamı çevrimiçi okunabilir, ayrıca PDF ve EPUB olarak indirilebilir. Ciltler Zenodo üzerinden DOI numaralarıyla arşivlenmiştir. Seri büyümeye devam etmektedir.

\vspace{0.6cm}
\noindent
\begin{minipage}{\textwidth}
\centering
\foreach \i/\ad in {1/Kumdan Bilgisayar, 2/Uzaydaki Bilgisayar, 3/Dediğimi Yap, 4/Eldenin Yolculuğu} {%
    \begin{minipage}[t]{0.235\textwidth}
        \centering
        \includegraphics[width=\linewidth]{gorseller/by-cilt\i.jpg}\\[3pt]
        {\scriptsize \ad}\\
        {\scriptsize\itshape \i.\ cilt}
    \end{minipage}\hfill
}
\end{minipage}

\vspace{0.8cm}
\noindent\rule{\textwidth}{0.5pt}
\vspace{0.3cm}

\noindent
{\large\bfseries \url{https://bilgeveyonga.oguzergin.net}}

\medskip
\noindent
Seri, bu kitapla aynı anlayışla, Creative Commons Atıf-GayriTicari-Türetilemez 4.0 lisansıyla açık erişime sunulmuştur. Öğretmenler ve veliler kitapları sınıfta ve evde ücretsiz kullanabilir.
